\documentclass[10pt,a4paper]{amsart}
\usepackage{amsmath,amssymb,amsthm,graphicx}
% Retained minimal preamble for the Stacks Project chapter varieties.tex.
% Verbatim excerpt of the chapter preamble: Xy-pic setup (native preamble lines 14-18), the native theorem
% declarations with the mathematical macros (native preamble lines 154-217), and the project own no-op
% definitions of the reference/slogan/history display environments (native wrapper.tex lines 1-3) that the
% counted statement uses for its slogan. Copied mechanically; nothing added, removed or redefined.
% For commutative diagrams we use Xy-pic
\usepackage[all]{xy}

% We use 2cell for 2-commutative diagrams.
\xyoption{2cell}
% Theorem environments.
%
\theoremstyle{plain}
\newtheorem{theorem}[subsection]{Theorem}
\newtheorem{proposition}[subsection]{Proposition}
\newtheorem{lemma}[subsection]{Lemma}

\theoremstyle{definition}
\newtheorem{definition}[subsection]{Definition}
\newtheorem{example}[subsection]{Example}
\newtheorem{exercise}[subsection]{Exercise}
\newtheorem{situation}[subsection]{Situation}

\theoremstyle{remark}
\newtheorem{remark}[subsection]{Remark}
\newtheorem{remarks}[subsection]{Remarks}

\numberwithin{equation}{subsection}

% Macros
%
\def\lim{\mathop{\mathrm{lim}}\nolimits}
\def\colim{\mathop{\mathrm{colim}}\nolimits}
\def\Spec{\mathop{\mathrm{Spec}}}
\def\Hom{\mathop{\mathrm{Hom}}\nolimits}
\def\Ext{\mathop{\mathrm{Ext}}\nolimits}
\def\SheafHom{\mathop{\mathcal{H}\!\mathit{om}}\nolimits}
\def\SheafExt{\mathop{\mathcal{E}\!\mathit{xt}}\nolimits}
\def\Sch{\mathit{Sch}}
\def\Mor{\mathop{\mathrm{Mor}}\nolimits}
\def\Ob{\mathop{\mathrm{Ob}}\nolimits}
\def\Sh{\mathop{\mathit{Sh}}\nolimits}
\def\NL{\mathop{N\!L}\nolimits}
\def\CH{\mathop{\mathrm{CH}}\nolimits}
\def\proetale{{pro\text{-}\acute{e}tale}}
\def\etale{{\acute{e}tale}}
\def\QCoh{\mathit{QCoh}}
\def\Ker{\mathop{\mathrm{Ker}}}
\def\Im{\mathop{\mathrm{Im}}}
\def\Coker{\mathop{\mathrm{Coker}}}
\def\Coim{\mathop{\mathrm{Coim}}}

% Boxtimes
%
\DeclareMathSymbol{\boxtimes}{\mathbin}{AMSa}{"02}

%
% Macros for moduli stacks/spaces
%
\def\QCohstack{\mathcal{QC}\!\mathit{oh}}
\def\Cohstack{\mathcal{C}\!\mathit{oh}}
\def\Spacesstack{\mathcal{S}\!\mathit{paces}}
\def\Quotfunctor{\mathrm{Quot}}
\def\Hilbfunctor{\mathrm{Hilb}}
\def\Curvesstack{\mathcal{C}\!\mathit{urves}}
\def\Polarizedstack{\mathcal{P}\!\mathit{olarized}}
\def\Complexesstack{\mathcal{C}\!\mathit{omplexes}}
% \Pic is the operator that assigns to X its picard group, usage \Pic(X)
% \Picardstack_{X/B} denotes the Picard stack of X over B
% \Picardfunctor_{X/B} denotes the Picard functor of X over B
\def\Pic{\mathop{\mathrm{Pic}}\nolimits}
\def\Picardstack{\mathcal{P}\!\mathit{ic}}
\def\Picardfunctor{\mathrm{Pic}}
\def\Deformationcategory{\mathcal{D}\!\mathit{ef}}
\newenvironment{reference}{}{}
\newenvironment{slogan}{}{}
\newenvironment{history}{}{}
\usepackage{cleveref}
\crefname{originallemma}{original lemma}{original lemmas}
\Crefname{originallemma}{Original lemma}{Original lemmas}
\makeatletter
\newlabel{lemma-geometrically-normal-in-codim-1}{{0C3P}{source}{}{sourcepointer-1}{}}
\newlabel{lemma-geometrically-normal-in-codim-1@cref}{{[originallemma][0][]0C3P}{[1][0][]source}{}{}{}}
\makeatother
\crefname{originallemma}{original lemma}{original lemmas}
\Crefname{originallemma}{Original lemma}{Original lemmas}
\makeatletter
\newlabel{lemma-pre-delta-invariant}{{0C3S}{source}{}{sourcepointer-2}{}}
\newlabel{lemma-pre-delta-invariant@cref}{{[originallemma][0][]0C3S}{[1][0][]source}{}{}{}}
\makeatother
\crefname{originallemma}{original lemma}{original lemmas}
\Crefname{originallemma}{Original lemma}{Original lemmas}
\makeatletter
\newlabel{lemma-pre-pre-delta-invariant}{{0C3R}{source}{}{sourcepointer-3}{}}
\newlabel{lemma-pre-pre-delta-invariant@cref}{{[originallemma][0][]0C3R}{[1][0][]source}{}{}{}}
\makeatother
\crefname{originalreference}{original reference}{original references}
\Crefname{originalreference}{Original reference}{Original references}
\makeatletter
\newlabel{algebra-lemma-pullback-module}{{02M1}{source}{}{sourcepointer-4}{}}
\newlabel{algebra-lemma-pullback-module@cref}{{[originalreference][0][]02M1}{[1][0][]source}{}{}{}}
\makeatother
\title{The delta-invariant at a point is unchanged by extension of the base field under the geometric normality hypotheses}
\author{The Stacks Project Authors}
\date{stacks-project:varieties.tex:local-snapshot-20261011; GNU FDL 1.2 or later}
\begin{document}
\maketitle
\noindent\textit{Editorial note.} All delivered mathematical text is the author's own, copied verbatim from the pinned Stacks Project chapter file varieties.tex; no statement or proof step was rewritten, reordered or removed.\par\medskip
\noindent\textit{Editorial note.} The counted span cites five cross-chapter labels, each recorded in assembly.external\_references with its canonical tag as a navigation pointer.\par\medskip
\noindent\textit{Editorial note.} The proof cites no bibliography item inside the counted span.\par\medskip
\noindent\textit{Allowed standard prerequisites (editorial).}\begin{itemize}
\item Geometrically normal in codimension one: assumptions and consequences. Stacks Project, Varieties, Lemma lemma-geometrically-normal-in-codim-1 (tag 0C3T). Supplies geometric reducedness of A and normality of A' tensor\_k K, which is what upgrades the inequality to an equality.
\item The pre-delta construction. Stacks Project, Varieties, Lemma lemma-pre-delta-invariant (tag 09NL). Provides the intermediate ring A' whose base change is compared on both sides.
\end{itemize}\medskip
\ifdefined\hypertarget\hypertarget{sourcepointer-1}{}\fi
\noindent\textit{Source reference.} Original lemma 0C3P: sources/stacks-chow/tags.txt, source lines 15199--15199. Cross-chapter standard prerequisite cited by the author inside the counted span; it is used as a pointer and its own proof is not part of this item.\par\medskip
\ifdefined\hypertarget\hypertarget{sourcepointer-2}{}\fi
\noindent\textit{Source reference.} Original lemma 0C3S: sources/stacks-chow/tags.txt, source lines 15202--15202. Cross-chapter standard prerequisite cited by the author inside the counted span; it is used as a pointer and its own proof is not part of this item.\par\medskip
\ifdefined\hypertarget\hypertarget{sourcepointer-3}{}\fi
\noindent\textit{Source reference.} Original lemma 0C3R: sources/stacks-chow/tags.txt, source lines 15201--15201. Cross-chapter standard prerequisite cited by the author inside the counted span; it is used as a pointer and its own proof is not part of this item.\par\medskip
\ifdefined\hypertarget\hypertarget{sourcepointer-4}{}\fi
\noindent\textit{Source reference.} Original reference 02M1: sources/stacks-chow/tags.txt, source lines 3428--3428. Cross-chapter standard prerequisite cited by the author inside the counted span; it is used as a pointer and its own proof is not part of this item.\par\medskip
\begin{lemma}
\label{lemma-delta-invariant-and-change-of-fields-better}
Let $k$ be a field. Let $X$ be a locally algebraic $k$-scheme.
Let $K/k$ be a field extension and set $Y = X_K$.
Let $y \in Y$ with image $x \in X$.
Assume assumptions (a), (b), (c) of
Lemma \ref{lemma-geometrically-normal-in-codim-1}
hold for $x \in X$ and that $\dim(\mathcal{O}_{Y, y}) = 1$.
Then the $\delta$-invariant
of $X$ at $x$ is $\delta$-invariant of $Y$ at $y$.
\end{lemma}

\begin{proof}
Set $A = \mathcal{O}_{X, x}$ and $B = \mathcal{O}_{Y, y}$.
By Lemma \ref{lemma-geometrically-normal-in-codim-1}
we see that $A$ is geometrically reduced.
Hence $B$ is a localization of $A \otimes_k K$.
Let $A \to A'$ be as in Lemma \ref{lemma-pre-delta-invariant}.
By Lemma \ref{lemma-geometrically-normal-in-codim-1}
we see that $A' \otimes_k K$ is normal.
Hence
$$
B' = B \otimes_{(A \otimes_k K)} (A' \otimes_k K)
$$
is normal, finite over $B$, and $B \to B'$ induces
an isomorphism on total rings of fractions. Namely, pick $f \in \mathfrak m_A$
satisfying (1) -- (6) of Lemma \ref{lemma-pre-pre-delta-invariant};
since $\dim(B) = 1$ we see that $f \in \mathfrak m_B$
playes the same role for $B$ and we see that $B_f = B'_f$ because $A_f = A'_f$.
It follows that $B \to B'$ is as in Lemma \ref{lemma-pre-delta-invariant}
for $B$. Thus we have to show that
$\text{length}_A(A'/A) =
\text{length}_B(B'/B) = \text{length}_B((A'/A) \otimes_A B)$.
Since $A \to B$ is flat (as a localization of $A \to A \otimes_k K$)
and since $\mathfrak m_B = \mathfrak m_A B$ (because
$B/\mathfrak m_A B$ is zero dimensional by the remarks above and
a localization of $K \otimes_k \kappa(x)$ which is reduced as
$\kappa(x)$ is separable over $k$) we conclude by
Algebra, Lemma \ref{algebra-lemma-pullback-module}.
\end{proof}

\end{document}
